% ============================== CHƯƠNG 2 ==============================
\chapter{Cơ sở lý thuyết}
\label{chap:background}

\section{Big Data và mô hình 5V}
\todo{Định nghĩa Big Data, 5V, so sánh với dữ liệu truyền thống~\cite{gandomi2015beyond}.}

\section{Kiến trúc Lakehouse và Medallion}
\todo{Data Lake vs Data Warehouse vs Lakehouse; tầng Bronze/Silver/Gold~\cite{armbrust2021lakehouse}.}

\section{Mô hình hoá kho dữ liệu: Kimball và Inmon}
\todo{Star schema, snowflake, fact/dimension, surrogate key; Role-playing, Junk và Degenerate dimension~\cite{kimball2013data}.}

\section{OLAP và các phép toán đa chiều}
\todo{Roll-up, drill-down, slice, dice, pivot; MOLAP/ROLAP; columnar storage.}

\section{Xử lý phân tán với Apache Spark}
\todo{RDD/DataFrame, lazy evaluation, DAG, Catalyst, broadcast join, AQE~\cite{zaharia2016spark}.}

\section{Ba cấp độ phân tích: Descriptive, Predictive, Prescriptive}
\todo{Định nghĩa và vị trí của từng trang dashboard trong chuỗi phân tích.}

\section{Học máy cho bài toán phân loại}
\subsection{Các thuật toán}
Logistic Regression, Decision Tree~\cite{breiman1984cart}, Random Forest~\cite{breiman2001random},
Gradient Boosting: XGBoost~\cite{chen2016xgboost} và LightGBM~\cite{ke2017lightgbm}.
\todo{Mô tả ngắn từng thuật toán.}

\subsection{Đánh giá mô hình}
Accuracy, Precision, Recall, F1, ROC-AUC, PR-AUC, Brier score; cross-validation theo nhóm;
out-of-time validation; permutation importance~\cite{breiman2001random}.

\subsection{Target leakage}
Rò rỉ nhãn xảy ra khi mô hình dùng thông tin chỉ có \emph{sau} thời điểm dự báo, dẫn tới
độ chính xác ảo khi huấn luyện nhưng vô dụng khi triển khai~\cite{kaufman2012leakage}.

\subsection{Sai số Bayes và trần lý thuyết}
Với phân phối $P(y\mid x)$, bộ phân loại Bayes $h^*(x)=\mathbb{1}[P(y=1\mid x)\ge 0.5]$ đạt sai số
nhỏ nhất có thể; không mô hình nào trên cùng tập đặc trưng có thể vượt qua một cách hệ thống~\cite{hastie2009elements}.
